% Generated draft: requires mathematical and source review.
\documentclass[11pt]{article}
\usepackage[margin=25mm]{geometry}
\usepackage{fontspec}
\setmainfont{Latin Modern Roman}
\usepackage{amsmath,amssymb,mathtools,unicode-math}
\setmathfont{Latin Modern Math}
\usepackage{enumitem,needspace,xurl,hyperref}
\setlength{\parindent}{0pt}
\setlength{\parskip}{5pt}
\setlength{\emergencystretch}{3em}
\newcommand{\N}{\mathbb{N}}
\newcommand{\Z}{\mathbb{Z}}
\newcommand{\Q}{\mathbb{Q}}
\newcommand{\R}{\mathbb{R}}
\newcommand{\C}{\mathbb{C}}
\newcommand{\F}{\mathbb{F}}
\newcommand{\A}{\mathbb{A}}
\newcommand{\PP}{\mathbb{P}}
\newcommand{\E}{\mathbb{E}}
\newcommand{\eps}{\varepsilon}
\newcommand{\dd}{\,\mathrm{d}}
\begin{document}
\section*{UnsolvedMath ID: 30004348}
\textbf{Source problem code:} catalogue entry
\subsection{Definitions and mathematical statement}
The two constructions named in the source are the extraction of \(L\)-functions from constant terms of Eisenstein series and the construction of \(L\)-functions by twisted doubling integrals. The covering groups under consideration are Brylinski--Deligne covering groups. The source does not specify the underlying field, reductive group, covering data, automorphic representations, Eisenstein series, or twisting data, and therefore does not define particular \(L\)-functions to compare.

\textbf{Comparison problem.} For Brylinski--Deligne covering groups, determine how the \(L\)-functions obtained from constant terms of Eisenstein series compare with the \(L\)-functions obtained from twisted doubling integrals.

\textbf{Unresolved qualifications.} This is a comparison question rather than a fully specified conjecture. The source does not state the range of groups and representations for which both constructions are available, how their input data should correspond, or what relation is sought. In particular, it does not specify whether the intended comparison concerns local factors, partial global \(L\)-functions, completed \(L\)-functions, normalizations, or functional equations. No equality or universal quantifier over such data can be supplied from this record.
\subsection{Short English statement}
The problem asks how two constructions of \(L\)-functions for Brylinski--Deligne covering groups are related: one using Eisenstein constant terms and the other using twisted doubling integrals. The supplied record leaves the input data and intended comparison unspecified, so a precise quantified conjecture requires further source clarification.
\subsection{Sources}
\begin{itemize}
\item Local UnsolvedMath registry, permanent ID \(30004348\), source problem code \texttt{catalogue entry}, title: Comparing Eisenstein and Twisted Doubling L Functions. The record belongs to the Oberwolfach Reports Open Problems collection; its proposed year is \(2019\), and no individual proposer is supplied.
\item Supplied report citation: \emph{Modular Forms (2020)}, Oberwolfach Reports, pp.\ 3529--3606, DOI \texttt{10.4171/owr/2019/57}. \url{https://doi.org/10.4171/owr/2019/57}. The background describes Oberwolfach Reports as documenting mathematical workshops and their open-problem discussions.
\item Supplied literature reference: Cai--Friedberg--Ginzburg--Kaplan, doubling constructions. \url{https://arxiv.org/abs/1710.00905}.
\item Supplied literature assessment, checked \texttt{2026-08-22}: the record is labeled open and described as an underspecified comparison question. The assessment states that twisted-doubling integral representations are developed for several classical-group \(L\)-functions, but that no general identification answering this extracted question is established. Its listed sources are the doubling-constructions reference and the OWR report above. This is supplied dated curation, not independently verified literature assessment; the source explicitly requests independent verification of time-sensitive status and cited claims.
\item The supplied statement status is \texttt{corrected\_verified}; its accompanying qualification says that only formula delimiters, notation, grammar, and typography were repaired, with mathematical content and scope preserved. This metadata does not establish mathematical correctness or resolve the formulation gaps.
\end{itemize}
\end{document}
